\section{Introducere}
\label{sec:intro}

În contextul tehnologic actual, caracterizat prin expansiunea accelerată a paradigmei IoT (\textit{Internet of Things}), arhitecturile sistemelor distribuite au suferit transformări fundamentale. Odată cu migrarea masivă a serviciilor către infrastructuri de tip Cloud-native, fie sub formă de SaaS (\textit{Software as a Service}), PaaS sau arhitecturi Serverless, complexitatea operațională a crescut exponențial. Aceste medii dinamice impun cerințe stricte privind observabilitatea, monitorizarea și controlul în timp real al resurselor.

Problematica monitorizării infrastructurilor moderne provine din limitările inherente ale arhitecturilor centralizate tradiționale. Deși eficiente la scară mică, sistemele bazate pe un coordonator unic care interoghează agenți pasivi devin rapid ineficiente în fața volumului de date generat de mii de containere sau microservicii efemere. Arhitecturile centralizate prezintă două vulnerabilități critice:

\begin{itemize}
    \item \textbf{Dependența de un Punct Unic de Eșec (SPOF):} În modelul clasic, starea globală a sistemului este agregată de un singur nod central. Această centralizare creează un \textit{Single Point of Failure}. Indisponibilitatea serverului central, fie din cauze hardware, fie software, duce la pierderea totală a vizibilității asupra întregii infrastructuri, lăsând administratorii fără date critice exact în momentele de criză.
    
    \item \textbf{Limitări de Scalabilitate și Latență:} Pe măsură ce numărul de noduri monitorizate ($N$) crește, pe nodul central întâmplă un proces de \textit{bottleneck}. Procesarea secvențială sau concurentă a mii de cereri de stare duce la saturarea lățimii de bandă și a resurselor CPU. Soluțiile convenționale implică costuri ridicate:
    \begin{itemize}
        \item \textit{Scalabilitate verticală:} Upgrade-ul hardware al serverului central este o soluție temporară și extrem de costisitoare.
        \item \textit{Scalabilitate orizontală:} Introducerea de proxy-uri sau sharding adaugă un nivel de complexitate în configurare și mentenanță care poate depăși beneficiile.
    \end{itemize}
\end{itemize}

Această lucrare propune și analizează o soluție alternativă, robustă și descentralizată: utilizarea unui protocol de tip \textbf{Gossip} bazat pe strategia \textbf{Dynamic Average Consensus} \cite{spanos2005dynamic}. Prin adoptarea unui model de comunicare simetric \textit{Push-Pull}, se demonstrează posibilitatea monitorizării în timp real a stării globale a sistemului (ex. încărcarea medie a procesorului) fără a utiliza un coordonator central, eliminând astfel riscurile structurale ale abordărilor clasice.